\documentclass[10pt]{article}
\input{../notes-style.tex}

\usepackage{multicol}
\usepackage{array}

\title{Ch.~5 Cheat Sheet\\\large Hash Tables}
\author{}
\date{}

\begin{document}
\maketitle
\vspace{-2em}

\begin{multicols}{2}

\subsection*{The big idea}
A hash table maps \textbf{key} $\to$ \textbf{bucket index} via a hash
function $h$. Insert/find/remove become array-index operations. Expected
cost: $O(1)$. Worst case: $O(n)$.

\subsection*{Hash function desiderata}
\begin{enumerate}
    \item \textbf{Deterministic.} Same key $\to$ same index.
    \item \textbf{Fast.} Ideally a handful of integer ops.
    \item \textbf{Uniform.} Keys spread across all buckets.
\end{enumerate}
Classic methods: \textbf{division} ($k \bmod m$, pick $m$ prime),
\textbf{multiplication} ($\lfloor m (k A \bmod 1) \rfloor$ for
$A \approx 0.618$), \textbf{mid-square} (square, take middle digits).

\subsection*{Load factor}
\[
\alpha = \frac{n}{\text{capacity}}
\]
\textbf{Chaining:} works up to $\alpha \approx 1$; cost grows with $\alpha$.\\
\textbf{Open addressing:} must keep $\alpha < 1$; performance degrades
sharply as $\alpha \to 1$. Rule of thumb: resize at $\alpha = 0.75$.

\subsection*{Chaining}
Each bucket = linked list. Collisions chain onto it.\\
\textbf{Insert:} hash $\to$ bucket $\to$ prepend node. $O(1)$.\\
\textbf{Find:} hash $\to$ bucket $\to$ walk chain. $O(1 + \alpha)$
expected.\\
\textbf{Remove:} same as find + unlink.\\
\textbf{Deletion:} easy -- just remove from list.

\subsection*{Linear probing}
On collision, try index $+1$, $+2$, $\ldots$ until empty.\\
Pros: great cache behavior (contiguous probes).\\
Cons: \textbf{primary clustering} -- once a run starts, it grows and
attracts more collisions.\\
Cost (expected, uniform hashing): insert $\approx
\frac{1}{2}(1 + \frac{1}{1-\alpha})$; find-miss $\approx
\frac{1}{2}(1 + \frac{1}{(1-\alpha)^2})$.

\subsection*{Quadratic probing}
Try index $+1^2, +2^2, +3^2, \ldots$.\\
Pros: avoids primary clustering.\\
Cons: \textbf{secondary clustering} -- keys with the same initial hash
probe the same sequence. Needs capacity chosen so the probe sequence
covers all slots (e.g., capacity = prime, $\alpha < 0.5$).

\subsection*{Double hashing}
Two hash functions: $h_1(k)$ for initial slot, $h_2(k)$ for step size.
Probe $h_1(k), h_1(k) + h_2(k), h_1(k) + 2 h_2(k), \ldots$.\\
Pros: breaks both clustering modes; closest to uniform hashing in
practice.\\
Cons: two hashes computed. $h_2$ must never return 0 and should be
coprime with capacity.

\columnbreak

\subsection*{Probe-table deletion: tombstones}
Open addressing can't just erase a slot -- that would break the probe
chain of some later key. Instead, mark the slot as a \textbf{tombstone}
(``was occupied, now empty''). Find skips past; insert may reuse.
Periodic full rehash is needed to clean them up.

\subsection*{Resizing (rehashing)}
Trigger: $\alpha$ exceeds a threshold (e.g., 0.75).\\
New capacity: usually $\approx 2\times$ old; pick the next prime if
method needs it.\\
\textbf{Every element must be rehashed} -- new indices come from new
modulus. Simple copy is wrong.\\
Cost: $O(n)$ single rehash; amortized $O(1)$ per insert over sequence.

\subsection*{Direct hashing (perfect hashing)}
If keys \emph{are} small integers in a dense known range, use them as
indices directly. No hash function, no collisions.\\
Use when: counting occurrences of ASCII chars (256-slot table),
counting digit frequencies, histogramming ages, \ldots.\\
Anti-use: sparse key space (10 employee IDs drawn from 9-digit SSNs
$\to$ $10^9$-slot table).

\subsection*{Comparison at a glance}
\begin{tabular}{@{}l|lll@{}}
\textbf{Strategy}   & \textbf{Insert} & \textbf{Find}   & \textbf{Notes}\\
\hline
Chaining            & $O(1)$          & $O(1+\alpha)$   & tolerates $\alpha \geq 1$\\
Linear probing      & $O(1)$ avg      & $O(1)$ avg      & primary cluster\\
Quadratic probing   & $O(1)$ avg      & $O(1)$ avg      & secondary cluster\\
Double hashing      & $O(1)$ avg      & $O(1)$ avg      & closest to ideal
\end{tabular}

\subsection*{Hash table vs.\ BST vs.\ sorted array}
\begin{tabular}{@{}l|lll@{}}
                      & \textbf{HT}       & \textbf{BST}      & \textbf{sorted arr}\\
\hline
find                  & $O(1)^*$          & $O(\log n)$       & $O(\log n)$\\
insert                & $O(1)^*$          & $O(\log n)$       & $O(n)$\\
min/max               & $O(n)$            & $O(\log n)$       & $O(1)$\\
sorted iter           & $O(n \log n)$     & $O(n)$            & $O(n)$\\
range query           & bad               & good              & good\\
memory overhead       & pointers + empty  & 2 ptrs/node       & none
\end{tabular}
$^*$ amortized / expected.

\subsection*{Crypto hashing (5.9) is a different beast}
Requirements: preimage resistance, second-preimage resistance, collision
resistance, \emph{slowness} (for passwords). MD5 and SHA-1 are broken
for collision resistance; use SHA-256+. For passwords use bcrypt /
scrypt / Argon2. Never roll your own.

\subsection*{Top gotchas}
\begin{itemize}
    \item Quoting hash lookup as $O(1)$ without saying \emph{expected}.
    \item Open addressing deletion by clearing the slot (breaks chains).
    \item Poor hash function + adversarial input = hash flooding DoS.
    \item Resizing that copies but doesn't rehash.
    \item Using \texttt{std::unordered\_map} when you need ordered
    iteration or range queries (use \texttt{std::map}).
\end{itemize}

\end{multicols}
\end{document}
